After the final evaluation, 23 follow-up analyses (P1 to P23) asked why augmentation did not help and what would. Each was added to the pre-registration (versions 4 to 21) before it ran, and every read of the test records is logged: 18 reads after the final evaluation. Later scores on the test records are therefore optimistic and do not replace the pre-registered results, so two analyses scored the best pipeline on 977 Icentia11k patients~\cite{tan2022icentia} that no analysis had used. Table~\ref{tab:posthoc} lists every analysis with its numbers, and Figure~\ref{fig:progress} follows the test score through the main changes.

\paragraph{The pre-registered answer holds.} With five seeds instead of three, no pre-registered verdict changed (P3). On two other databases, St Petersburg INCART and the MIT-BIH Supraventricular Arrhythmia Database~\cite{goldberger2000}, generative augmentation was detectably better than no augmentation in 1 and worse in 1 of 8 comparisons (P2). With beats split at random, so that the same patients are in training and test, the same cnn reached 0.777 to 0.837 against 0.088 across patients (P4): much of what such splits measure is how well a model recognises patients it has seen.

\paragraph{The patient limits the model more than the class balance.} Giving the cnn each beat relative to the patient's own median beat (P6, Figure~\ref{fig:relative}), with nothing else changed, raised its test score without augmentation from 0.088 to 0.371 (P7), more than any rebalancing method did. The synthetic S beats came from the same few patients as the real ones: 0.61 of them were nearest to a beat of record 209, against 0.60 of the real training S beats (P8, Figure~\ref{fig:synthetic-patients}), so the generators added beats but not patients. Methods that act on the patient were then tested: moving the deviation of each rare beat onto other patients' normal beats (P9), normalising each patient inside the network (P10), and adding the INCART and SVDB patients (P11) or 400 Icentia11k patients (P14) to training. Of their 10 test comparisons, none changed detectably; at this scale, adding patients did not by itself raise the score.

\paragraph{Timing identifies S beats, as it does for a cardiologist.} An S beat is recognised mainly by coming early for the patient's own rhythm (Section~\ref{sec:medical}). In P13, dividing the RR intervals by the recording's mean RR interval, as published methods do, did not detectably change the cnn's test score in any of the 5 conditions (without augmentation 0.147 to 0.216, S F1 0.294 to 0.399); with Huang et al.'s rule on top, which uses the same feature directly, S F1 reached 0.531 to 0.578 (sensitivity 0.920 to 0.955) and the cVAE's rare-class macro-F1 0.523 (descriptive, not tested). Figure~\ref{fig:rr-reference} shows what each reference makes of the training-side recordings. In P16, undoing the rule's S calls inside atrial fibrillation, found by a detector that reads only beat times, raised rare-class macro-F1 in all 6 conditions (by 0.018 to 0.034; every 95\% interval above zero; detectable after Holm correction in 2), by removing false S calls: the cVAE went from 0.523 to 0.557 (S F1 0.572 to 0.640, F F1 0.474).

\paragraph{Better generators and a larger classifier did not change the answer.} A class-conditional diffusion model (P15), the same model trained five times longer (P17), the physics loss of SimGANs, which scores generated beats with an ECG simulator (P21), and a physics loss fitted to each generated beat (P23) were each compared on the test records with the generator they replace or with no augmentation; of these 8 comparisons, none changed detectably. A ResNet with 6.7 times the cnn's parameters (P22) raised supraventricular F1 on the unused patients in 2 of 3 comparisons; of its 3 comparisons on the test records, none changed detectably.

\paragraph{The best pipeline, and how far to trust it.} In P18, averaging the class probabilities of ten independently trained runs per condition changed rare-class macro-F1 by $-$0.004 to +0.099 against P16's single runs (after Holm correction: raised it in 1 of 6 conditions); the cVAE's ensemble reached 0.587, the highest test score of the study (S F1 0.647, F F1 0.527), against 0.557 for its single runs (not a detectable change, Holm-adjusted $p=1.000$). The mean of its S and F F1 is close to the best published inter-patient result we could check, 0.582 for Zhu et al.~\cite{zhu2022} (Section~\ref{sec:all-results}). The pipeline was designed after the test records had been read many times, so this score is optimistic. In P19, on 977 Icentia11k patients that no analysis had used, the same pipeline, unchanged, reached S F1 0.294 (95\% CI 0.237 to 0.350) and V F1 0.648 for the cVAE's ensemble with the stage and the veto, and S F1 0.348 for the ensemble alone. Icentia11k has no fusion label, so only S and V can be scored there. In P20, choosing the stage and the veto on 400 other Icentia11k patients picked \emph{stage 0.70, veto 0.80} for the cVAE, which took its S F1 on P19's patients from 0.294 to 0.414 (improved, Holm-adjusted $p=0.001$); on the test records its rare-class macro-F1 went from 0.587 with P18's setting to 0.459. On new patients, then, the pipeline's supraventricular F1 (0.294 and 0.414) stays well below the 0.647 measured on the test records.
